\section{Repository reuse and Lean formalization}
\label{sec:lean}

\subsection{The inspected foundation}

The ProveIt repository was inspected at the fixed commit
\[
 \texttt{\small 8c135f6f2a9d041cde785a1a632c44ad98065b20}.
\]
The snapshot's toolchain and Mathlib dependency are both pinned to
version \(4.32.0\). The relevant source directory is
\begin{center}
\code{Logic/PrimitiveRecursiveArithmetic/ConciseIPRA/Lean/ConciseIPRA/}.
\end{center}
The following declarations were read in source, rather than inferred
from file names or an informal summary:
\begin{itemize}
\item \code{PR.lean}: the types \code{PR} and \code{PTm}, the axiom
family \code{PRAAx}, and the rule \code{PRA.qfInd}.
\item \code{PRAToScalar.lean}: \code{FromPRA.code},
\code{FromPRA.interp}, and proofs translating every arithmetic axiom,
including the quantifier-free induction axiom.
\item \code{BiInterpretation.lean}: \code{biInterp},
\code{pra\_iff\_teq}, \code{teq\_iff\_pra}, and the round-trip theorems.
\item \code{NativeBiInterpretation.lean}: \code{NFm.ofPRA},
\code{pra\_iff\_native}, and \code{native\_iff\_pra}.
\end{itemize}
These sources establish the existing syntactic translation interfaces
\cite{proveit}. They do not already provide the finite sieve proved here.
The foundation was not rebuilt locally in producing this package.

The inspected draft
\code{docs/drafts/prime\_reciprocal\_ipra\_2026-11-15.zip}
contains useful bounded-factorization and harmonic-sum methodology
\cite{prime-draft}. Its Lean scaffold contains explicit unproved
placeholders. It is therefore background for the proof design, not an
importable completed number-theory proof. The present argument proves
all of its required arithmetic facts independently.

\subsection{Two distinct formalization goals}

There are two layers to complete.

\paragraph{Finite mathematical layer.}
Prove the rational finite theorems in Lean, using computable finite lists
or finite sets. The primary targets are the quadratic identity, the
harmonic comparison, the dyadic count, and the finite tail bound.
Lean's rational numbers are convenient for checking these identities.
No infinite sums or real analysis are needed.

\paragraph{Object-theory layer.}
Build actual primitive recursive syntax codes for the displayed numerical
functions and a derivation in the deeply embedded theory \code{PRAAx}.
The final target has the form
\[
 []\ \vdash_{\mathsf{PRAAx}}\ \mathrm{BrunFormula}.
\]
A Lean theorem asserting the ordinary truth of
\eqref{eq:ipra-target} is valuable, but it does not on its own establish
that iPRA proves it. The additional syntactic derivation is essential.

\subsection{Recommended modules and interfaces}

\renewcommand{\arraystretch}{1.12}
\begin{longtable}{@{}p{0.30\textwidth}p{0.63\textwidth}@{}}
\toprule
Proposed module & Completion target\\
\midrule
\endhead
\code{Brun/FiniteArithmetic} &
Prove correctness of the bounded prime test; construct the factor trace,
squarefree support, divisor list and rational codes.\\
\code{Brun/Residues} &
Prove the explicit CRT reconstruction, \(\rho(d)=\prod_{p\mid d}\rho(p)\),
and the rational error bound for \(A_d(X)\).\\
\code{Brun/SelbergIdentity} &
Define \(G,W,\lambda\); prove \(\lambda_1=1\),
\(|\lambda_d|\rho(d)\le d\), the transform, and \(Q=1/G\).
Then derive \(T(X)\le z+X/G+z^4\).\\
\code{Brun/DenominatorLower} &
Construct the injection from bounded pairs into prime exponent boxes;
prove finite product factorization and \(H_s^2\le3G(s^2)\).\\
\code{Brun/DyadicCount} &
Prove the harmonic and polynomial inequalities and instantiate
\(t=k\mathbin{\mathrm{div}}10\), obtaining \(k^2T(2^k)\le2048\,2^k\).\\
\code{Brun/Tail} &
Construct the bounded shell index, bound each shell, telescope the
finite sum, and prove \eqref{eq:main-tail}.\\
\code{Brun/Cauchy} &
Finish the supplied conditional assembly, and prove equivalence with
the factorial-denominator integer statement.\\
\code{Brun/IPRA/Codes} &
Construct explicit \code{PR} syntax trees and defining-equation
derivations for all operations used by the proof.\\
\code{Brun/IPRA/Derivation} &
Internalize the lemmas using \code{PRA.qfInd}, with a quantifier-free
syntactic witness at every arithmetic induction, and derive the closed
formula.\\
\code{Brun/IPRA/Native} &
Apply the existing faithful interpretations to the completed
\code{PRAAx} derivation.\\
\bottomrule
\end{longtable}

The most substantial new tasks are the finite prime-support injection and
the syntactic arithmetic library. The number-theoretic bottleneck is
specified by Theorem~\ref{thm:harmonic}; no unknown analytic estimate is
left for the formalization to discover.

\subsection{A concrete semantic specification}

The supplied \code{lean/BrunSpecification.lean} defines:
\begin{itemize}
\item \code{primeB} and \code{twinB} by bounded list scans;
\item \code{brunPartial}, with the strict endpoint \(p<N\);
\item \code{sieveConstant}, \code{exponentIndex}, and \code{cutoff},
with constants exactly as in this article;
\item \code{tailNumerator} and \code{ArithmeticBrunClaim},
which encode \eqref{eq:integer-tail} and \eqref{eq:ipra-target};
\item \code{DyadicTailClaim} and \code{BrunCauchyClaim};
\item candidate conditional assembly proofs showing how a proof of the
finite tail bound yields the stated Cauchy estimate.
\end{itemize}
There are no newly declared axioms or \code{sorry} placeholders in that
file. The missing sieve theorem is a named proposition supplied as an
\emph{explicit hypothesis} to the assembly lemma; it is not asserted as
proved. The source has not been compiled because a Lean/Lake executable
was unavailable in the authoring environment.

The crucial assembly principle is elementary. Suppose every \(a(n)\)
beyond a cutoff lies in the same interval
\([a(N),a(N)+\epsilon]\). Then any two such values differ in absolute value
by at most \(\epsilon\), not \(2\epsilon\). This explains why the tail
bound produces the stated modulus without an additional factor of two.

\subsection{Building the object-language formula}

The existing \code{PR} syntax is generated by zero, successor,
projections, composition and primitive recursion. Define within this syntax
the fixed operations needed for the proof: comparison, quotient,
remainder, exponentiation, factorial, bounded folds, twin testing,
the cutoff, and the integer tail numerator.

For each function, prove its computation equations in \code{PRAAx}.
For each finite lemma, prove the corresponding coded statement by the
specific invariant in Section~\ref{sec:ipra}. Algebraic normalization can
be automated by a proof-producing macro that emits uses of the semiring
and rational-code lemmas. A computation that merely reports a true
Boolean is not a substitute for a uniform derivation.

With variables ordered as \(q,U,V\), form a primitive recursive Boolean
\(b(q,U,V)\) that is \(1\) precisely when the antecedent of
\eqref{eq:ipra-target} is false or its conclusion is true.
Both tests are decidable. The equality \(b(q,U,V)=1\) is an atomic
formula \code{brunMatrix : Fm PTm 3}; its \code{Fm.QF} evidence is immediate.
Universally close the three variables to obtain
\code{brunFormula : Fm PTm 0}.

The final task is to construct a term of
\[
 \code{Prf}\ \mathsf{PRAAx}\ []\ \mathrm{brunFormula},
\]
using the repository's exact notation and parameter order.
All arithmetic induction applications must go through
\code{PRA.qfInd} or a derived macro that supplies the required
\code{Fm.QF} witness.
Structural induction in Lean on a syntax tree or on a derivation is
metatheory; it does not grant unrestricted numerical induction in the
object theory.

\subsection{Transfer and verification gates}

After a derivation \code{brunPRA} is constructed, its scalar translation
is obtained through \code{FromPRA.interp.sound}; the native translation is
available through \code{pra\_iff\_native}. Thus there is no additional
bi-interpretability theorem to prove for this number-theory application.

Completion should be checked in the pinned ProveIt environment:
\begin{enumerate}
\item Compile the finite arithmetic modules and the supplied specification.
\item Compile the \code{PR} codes and the actual
\code{PRAAx} derivation.
\item Compile the scalar and native corollaries.
\item Inspect \code{\#print axioms} for the final declarations and verify
that there is no \code{sorryAx} or assumption standing for Brun's theorem.
\item Check that the final result is a derivation in \code{PRAAx}, not
only a theorem about standard natural numbers.
\item Review the \code{Fm.QF} witness at every object-level induction.
\end{enumerate}
The accompanying \code{FORMALIZATION\_PLAN.md} gives the same route in a
form convenient for implementation. No successful Lean build is claimed
for this package.
